\section{Experimental Evaluation}\label{sec:experiments}

%\todo{put ratios in macros and update throughout}
We evaluate three succinct tries optimized using \csquared (\csquared-FST,
\csquared-CoCo, and \csquared-Marisa) and compare them with \numtries state-of-the-art
tries: FST \cite{surf}, Marisa \cite{marisa}, CoCo-trie \cite{coco_0, coco_1},
PDT \cite{pdt}, ART \cite{art}, C-ART \cite{hybrid_index}, \newmymarginpar{Meta
  4, R2 O2}\newrev{c-trie++~\cite{tsuruta2022c}, and z-fast
  trie~\cite{belazzougui2010dynamic}}. The first four are \bv-based succinct
tries introduced in \secref{preliminaries}; ART is a well-established
pointer-based dynamic trie, C-ART is the compact and static version of ART, and
\newrev{c-trie++/z-fast trie are high-performance dynamic tries}. We compare all
data structures on query performance, space usage, and build time.

\para{Summary} The \csquared tries significantly outperform the original
versions in both performance and space efficiency on most datasets.  On average,
\csquared-FST, \csquared-CoCo, and \csquared-Marisa improve query
performance/space usage by \csqfstspeedup/\csqfstspacesaving,
\csqcocospeedup/\csqcocospacesaving, and \\
\csqmarisaspeedup/\csqmarisaspacesaving, respectively (at equal recursion).

These improvements stem from improved locality and fewer cache misses in both
the index and the data. Applying \bvdesign (the cache-conscious bitvector
redesign in~\secref{bv-design}) reduces cache misses per query by $1.51\times$
for FST and $1.26\times$ for Marisa on large datasets. The \bv redesign adds at
most $4\%$ space overhead. At equal recursion, \compressionscheme (the
unary-path compression scheme in~\secref{unary-path}) improves query
performance, space usage, or both over \bvdesign alone.

\subsection{Experimental Settings}\label{sec:experimental-setup}


\begin{table}[t]
  \caption{Datasets used in evaluation. All dataset sizes are expressed in
    MB. LCP means length of longest common prefix.  \textit{Size*} is the size
    of the prefix-only version of each dataset.}
  \resizebox{\columnwidth}{!}{
 \newrev{  \begin{tabular}{l r r r r r l}
             \toprule

    \textit{Name} & \textit{Size} & \textit{Size*} & \textit{\#Keys}
    &\makecell{\textit{Avg} \\ \textit{len}} &\makecell{\textit{Avg} \\ \textit{LCP}} & \textit{Description}\\
    \midrule

    \wordsdata~\cite{words-dataset}
    & 5
    & 4
    & 4.67E5
    & 9
    & 6
    & \small{English words}\\

    \urldata~\cite{web-crawl}
    & 37
    & 17
    & 1.77E6
    & 21
    & 7
    & \small{UK private domains}\\

    \dnadata~\cite{text-collection}
    & 101
    & 44
    & 3.32E6
    & 31
    & 11
    & \small{DNA 31-mers}\\

    \xmldata~\cite{text-collection}
    & 117
    & 74
    & 2.15E6
    & 56
    & 33
    & \small{dblp XML dump} \\

    \logdata~\cite{log-dataset}
    & 586
    & 258
    & 4.45E6
    & 137
    & 54
    & \small{server access logs}\\

    \wikidata~\cite{wiki-dataset}
    & 359
    & 236
    & 1.75E7
    & 21
    & 11
    & \small{Wikipedia titles}\\
    \bottomrule
  \end{tabular}}
}
  \label{tab:datasets}
\end{table}


%%% Local Variables:
%%% mode: latex
%%% TeX-master: "../sample-sigconf"
%%% End:


\para{Hardware} \newrev{We ran the experiments on a server with an AMD EPYC 9555
  64-Core Processor running at 3.2GHz. The server has 1007 GB of memory, a 3 MiB
  L1 data cache, a 2 MiB L1 instruction cache, a 64 MiB L2 cache, and a 256 MiB
  L3 cache.} We ran the experiments with one thread.


\para{Implementation} We implement the three \csquared tries (\csquared-FST,
\csquared-CoCo and \csquared-Marisa) using C++17. We use the compressed string
pool implementation of PDT~\cite{pdt-code} \new{and the publicly-available
  implementation~\cite{fsst-code} of FSST~\cite{fsst} for the \repair and \fsst
  tail containers, respectively.} We also \rev{use} the sdsl~\cite{sds_overview,
  sdsl_lib}, ds2i~\cite{ds2i} and sux~\cite{rank_select_0} libraries for
\csquared-CoCo as in the original CoCo-trie~\cite{coco-code}.
%\new{Finally, we use the optimized build routine from CoCo' for \csquared-CoCo.}

For CoCo-trie~\cite{coco-code}, Marisa~\cite{marisa-code}, PDT~\cite{pdt-code}
and C-ART~\cite{cart-code}, we use their original implementations. We adopt the
ART implementation in the C-ART codebase. For FST, we use an optimized
third-party implementation~\cite{fst-code} which compacts suffixes into a
contiguous array, as this allows for fairer comparison with many unary suffix
paths. \newmymarginpar{Meta 4, R2 O2}\newrev{Finally, we also evaluate
  c-trie++~\cite{tsuruta2022c} and z-fast trie~\cite{belazzougui2010dynamic},
  two state-of-the-art dynamic tries. }

\newmymarginpar{Meta 3, R3 O2}\newnew{All tested systems support
  existence queries, but out of the tested \csquared tries, only FST supports an
  iterator for range queries. Therefore, we evaluate FST versus \csquared-FST on
  range queries. In each experiment, we generate range queries with a start key
  and a length, or number of subsequent keys to return beginning from the start
  key. We vary the query lengths \texttt{k} between experiments and test
  $\texttt{k} = 1, 10, \ldots, 10000$.}

The ART and C-ART implementations we use are designed as filters but not indexes
and may return false positives. To support lossless queries, the original
dataset must be replicated and may need to be checked when the filter returns
positive results. We report space costs without the replicated dataset; query
times include false-positive checks.

Finally, as detailed in~\secref{coco-prime}, we introduce CoCo', a version of
the CoCo-trie with an optimized build routine but the same \bv design, because
the open-source CoCo-trie does not build on large datasets.

All code is compiled with g++ 10.3.0 at -O3.

All implementations used in our experiments are publicly available at
\url{https://github.com/alexztc/C2}.

\para{Parameter settings} For FST, we use the default setting with $R =
64$. \rev{For CoCo-trie, CoCo', and \csquared-CoCo, we set $\alpha = 5\%$ for a
  good space-time balance.}  We do not enable LOUDS-Dense for \csquared-FST as
it provides no benefit. \new{We use Marisa's default cache size for both Marisa
  and \csquared-Marisa, which is 1/512 of the key count. We set $\epsilon = .1$
  for all \csquared tries.}

\para{Datasets} \tabref{datasets} details the \numdatasets datasets in our
evaluation.  These datasets span diverse sources, key lengths, prefix depths,
and alphabet sizes.

\rev{Furthermore, following the methodology from the CoCo-trie
  paper~\cite{coco_1}, we generate prefix-only versions of each dataset (denoted
  by \texttt{dataset*}) because CoCo originally used only prefix-only
  datasets. We compare the open-source CoCo implementation with our CoCo'
  implementation on the prefix-only datasets.}

\para{Experimental setup} We set up experiments as follows: First, each dataset
is sorted and deduplicated. Second, we build the trie under test on the dataset
and record the build time and memory usage. Finally, we query the trie using all
keys in the dataset in random order and compute the average query latency. We
omit negative-query latencies, as they are generally proportional to positive
workloads~\cite{coco_1}. We also measure the cache performance of queries in
different tries on the \logdata and \wikidata datasets by recording the number
of cache misses using \texttt{perf}. \newrev{All times are the average of three
  trials after one warm-up trial.}

\para{Notation} We use Marisa-$i$ to denote Marisa with $i$ levels of
recursion. Furthermore, for a given trie X, we use \bvdesign-X to denote X
optimized with \bvdesign (\secref{bv-design}). We use \csquared-X to denote X
with both \bvdesign and \compressionscheme. For example, \bvdesign-FST is
\bvdesign-optimized FST with sorted tail container, whereas \csquared-Marisa-1
is \bvdesign-optimized Marisa with \fsst tail container and one level of
recursion.


\subsection{Building \csquared-CoCo on larger datasets}\label{sec:coco-prime}
% %% update words and url
% \begin{table}[t]
%     \caption{
%   Comparison of CoCo and CoCo$'$ (\queryshort, in ms for total queries) and space usage (\sizeshort,
%   in \% of original dataset size) on the
%   prefix-only datasets used in the original CoCo-trie
%   evaluation.}
%       \begin{center}
%           \newcolumntype{P}{>{\raggedleft\arraybackslash}p{.24in}}
%           \resizebox{\columnwidth}{!}{
%            \newrev{  \begin{tabular}{l P P P P P P P P }
%               \toprule

%                & \multicolumn{2}{c}{\wordsstar}
%                & \multicolumn{2}{c}{\urlstar}
%                & \multicolumn{2}{c}{\dnastar}
%                & \multicolumn{2}{c}{\xmlstar}\\

%               \cmidrule(lr){2-3}
%               \cmidrule(lr){4-5}
%               \cmidrule(lr){6-7}
%               \cmidrule(lr){8-9}

%              \textit{Trie} &
%              \textit{\queryshort} & \textit{\sizeshort} &
%              \textit{\queryshort} & \textit{\sizeshort} &
%              \textit{\queryshort} & \textit{\sizeshort} &
%              \textit{\queryshort} & \textit{\sizeshort} \\

%               \midrule

%               % CoCo
%               % & 2278 & 62.7\%
%               % & 4216 & 42.4\%
%               % & 654 & 24.0\%
%               % & 1056 & 47.5\% \\

%               % CoCo$'$
%               % & 1536 & 54.8\%
%               % & 2929 & 33.5\%
%               % & 742 & 40.1\%
%               % & 1020 & 29.6\% \\

%                 CoCo
%                 & 2278 & 62.7\%
%                 & 4216 & 42.4\%
%                 & 3052 & 24.0\%
%                 & 7696 & 47.5\% \\

%                 CoCo$'$
%                 & 1536 & 54.8\%
%                 & 2929 & 33.5\%
%                 & 2344 & 73.2\%
%                 & 3846 & 40.6\% \\

%               \bottomrule
%           \end{tabular}}
%           }
%       \end{center}

%   \label{tab:coco-vs-coco-prime}
%   \end{table}

% %%% Local Variables:
% %%% mode: latex
% %%% TeX-master: "../main"
% %%% End:



  \begin{table}[t]
      \caption{
    Comparison of CoCo and CoCo$'$ (\queryshort, in ns per query) and space usage (\sizeshort,
    in \% of original dataset size) on the
    prefix-only datasets used in the original CoCo-trie
    evaluation.}
        \begin{center}
            \newcolumntype{P}{>{\raggedleft\arraybackslash}p{.24in}}
            \resizebox{\columnwidth}{!}{
             \newrev{  \begin{tabular}{l P P P P P P P P }
                \toprule

                 & \multicolumn{2}{c}{\wordsstar}
                 & \multicolumn{2}{c}{\urlstar}
                 & \multicolumn{2}{c}{\dnastar}
                 & \multicolumn{2}{c}{\xmlstar}\\

                \cmidrule(lr){2-3}
                \cmidrule(lr){4-5}
                \cmidrule(lr){6-7}
                \cmidrule(lr){8-9}

               \textit{Trie} &
               \textit{\queryshort} & \textit{\sizeshort} &
               \textit{\queryshort} & \textit{\sizeshort} &
               \textit{\queryshort} & \textit{\sizeshort} &
               \textit{\queryshort} & \textit{\sizeshort} \\

                \midrule

                  CoCo
                  & 315 & 62.7\%
                  & 742 & 42.4\%
                  & 655 & 24.0\%
                  & 1039 & 47.5\% \\

                  CoCo$'$
                  & 380 & 54.8\%
                  & 811 & 32.3\%
                  & 743 & 40.1\%
                  & 1011 & 29.6\% \\

                \bottomrule
            \end{tabular}}
            }
        \end{center}

    \label{tab:coco-vs-coco-prime}
    \end{table}
Since the original CoCo-trie implementation
runs out of memory when building the larger prefix-only datasets, we implement
an optimized build routine underneath \csquared-CoCo to support all datasets.
The original CoCo is built from a pointer-based uncompacted trie, which not only
consumes excessive memory but also incurs significant cache misses when
traversing each node's descendants at different levels, a key operation of
CoCo-trie's optimization process. We optimize the builder by representing the
uncompacted trie as \csquared-FST, which significantly reduces cache misses
because in LOUDS-Sparse representation each node's descendants at the same level
fall into contiguous memory. Our implementation reduces the build time of
\csquared-CoCo by up to $14\times$ and uses orders of magnitude less memory
during build than the original CoCo-trie.

We introduce \defn{CoCo'}: \csquared-CoCo's build routine with the original
CoCo-trie \bv, to isolate \bv differences. % on larger datasets.

\mymarginpar{R2 D8, R2 D11}\rev{We use CoCo' as the CoCo-trie baseline for
  larger datasets. As shown in~\tabref{coco-vs-coco-prime}, on the prefix-only
  datasets that the original CoCo-trie can build on, CoCo-trie and CoCo' have
  similar query latencies (within \newrev{100} ns). However, the original CoCo-trie suffers from poor
  performance and space usage on full datasets because it does not handle unary
  suffix paths. For example, on \urldata, a relatively small dataset, the
  original CoCo-trie implementation requires more than 10 minutes and 60 GB of
  memory to build. Furthermore, it takes over 2000 ns/query and occupies nearly
  half as much memory as the original dataset, both the worst among all tested
  tries.}


% \begin{table}[t]
%     \caption{Build and query time (normalized to \csquared-FST, lower is better)
%     for different LOUDS-Sparse/Dense configurations for FST. FST-Sparse/Hybrid
%     use LOUDS-Sparse/LOUDS-Dense and LOUDS-Sparse with a sorted tail
%     container. \csquared-FST uses the same bitvectors with the FSST tail
%     container.}
%   \centering \setlength{\tabcolsep}{5pt} \newnew{\begin{tabular}{lrrrr} 
  
%   \toprule
%                                                    & \multicolumn{2}{c}{\textit{words} (0.47M keys)} & \multicolumn{2}{c}{\textit{log} (4.45M keys)} \\
%                                                    \cmidrule(lr){2-3}\cmidrule(lr){4-5}
%                                                    & \textit{Build} & \textit{Query} & \textit{Build} & \textit{Query} \\
%                                                    \midrule
%                                                    FST-Sparse & 7.73$\times$ & 3.50$\times$ & 2.30$\times$ & 2.73$\times$ \\
%                                                    FST-Hybrid & 7.02$\times$ & 2.85$\times$ & 2.24$\times$ & 2.57$\times$ \\
%                                                    C\textsuperscript{2}-FST-Sparse & 5.31$\times$ & 1.52$\times$ & 1.53$\times$ & 1.82$\times$ \\
%                                                    C\textsuperscript{2}-FST-Hybrid & 5.29$\times$ & 1.45$\times$ & 1.55$\times$ & 1.85$\times$ \\
%                                                    \midrule
%                                                    C\textsuperscript{2}-FST$$ & 1.00$\times$ & 1.00$\times$ & 1.00$\times$ & 1.00$\times$ \\
%                                                    \bottomrule
%     \end{tabular}}
%   \label{tab:fst-ablation-rel}
% \end{table}

% %%% Local Variables:
% %%% mode: latex
% %%% TeX-master: "../main"
% %%% End:


\begin{table}[t]
\caption{
Normalized build and query time (relative to \csquared-FST; lower is better)
for different LOUDS-Sparse/Dense configurations of FST.
FST-Sparse uses LOUDS-Sparse/LOUDS-Dense, while FST-Hybrid uses
LOUDS-Sparse, both with a sorted tail container.
\csquared-FST uses the same bitvectors with the FSST tail container.
}
\centering
\setlength{\tabcolsep}{5pt}
\begin{tabular}{lrrrr}
\toprule
& \multicolumn{2}{c}{\textit{words} (0.47M keys)}
& \multicolumn{2}{c}{\textit{log} (4.45M keys)} \\
\cmidrule(lr){2-3}
\cmidrule(lr){4-5}
& \textit{Build}
& \textit{Query}
& \textit{Build}
& \textit{Query} \\
\midrule
FST-Sparse
    & 7.73$\times$
    & 3.50$\times$
    & 2.30$\times$
    & 2.73$\times$ \\

FST-Hybrid
    & 7.02$\times$
    & 2.85$\times$
    & 2.24$\times$
    & 2.57$\times$ \\

\csquared-FST-Sparse
    & 5.31$\times$
    & 1.52$\times$
    & 1.53$\times$
    & 1.82$\times$ \\

\csquared-FST-Hybrid
    & 5.29$\times$
    & 1.45$\times$
    & 1.55$\times$
    & 1.85$\times$ \\
\iffalse
\midrule

\csquared-FST
    & 1.00$\times$
    & 1.00$\times$
    & 1.00$\times$
    & 1.00$\times$ \\
\fi
\bottomrule
\end{tabular}
\label{tab:fst-ablation-rel}
\end{table}


%%% New organization

\begin{table*}[t]
 
\caption{
%\todo{add more recursion levels for the other to get the pareto frontier?}
Build time,
query latency,
and space usage on the datasets in~\tabref{datasets}.
%Bolded numbers are the best in each column.
The prefixes
        \bvdesign/\csquared indicate that the
        \bvdesign/\bvdesign+\compressionscheme optimizations are applied.  We
        fix the tail container to be FSST with the \compressionscheme
        optimization, but we report results with the sorted tail container for
        the ablation study on the \bvdesign optimization. Marisa-1 means the
        Marisa trie with 1 recursion level.  For Marisa, we show all data points
        before \compressionscheme stops recursion.  For all other tries, the
        recursion level was set to 0 for a fair
        comparison. For each column, green cells highlight the best value and blue cells highlight the second-best value.
}
    \begin{center}
        \newcolumntype{P}{>{\raggedleft\arraybackslash}p{.24in}}
        \resizebox{\textwidth}{!}{
\newrev{          \begin{tabular}{r P P P P P P P P P P P P P P P P P P}
 
            \toprule
 
             & \multicolumn{6}{c}{\textit{\buildtimeshort}   \textit{(ns/key)}}
             & \multicolumn{6}{c}{\textit{\queryshort} \textit{ (ns/key)}}
             & \multicolumn{6}{c}{\textit{\sizeshort} \textit{(\% of original size)}} \\
 
             \cmidrule(lr){2-7}
             \cmidrule(lr){8-13}
             \cmidrule(lr){14-19}
 
           \textit{Trie} &
           \textit{\wordsdata} & \textit{\urldata} & \textit{\dnadata} &
           \textit{\xmldata} & \textit{\wikidata} & \textit{\logdata} &
           \textit{\wordsdata} & \textit{\urldata} & \textit{\dnadata} &
           \textit{\xmldata} & \textit{\wikidata} & \textit{\logdata} &
           \textit{\wordsdata} & \textit{\urldata} & \textit{\dnadata} &
           \textit{\xmldata} & \textit{\wikidata} & \textit{\logdata} \\
 
            \midrule
 
FST
& 309&691&802&1157&716&2139
& 403&610&844&1541&1424&4003
& 45.5\%&35.5\%&74.3\%&39.2\%&37.9\%&39.5\%\\
 
\bvdesign-FST
& \cellcolor{blue!20}90&1043&1148&897&762&2173
& 228&404&562&835&1242&2197
& 40.4\%&39.8\%&75.9\%&39.7\%&38.4\%&40.4\%\\
 
\csquared-FST
& 101&1068&1114&894&729&1821
& 227&404&561&842&1258&2154
& 40.4\%&36.0\%&43.8\%&28.8\%&37.7\%&24.4\%\\
 
\midrule
 
% CoCo'
% & 790&1639&1875&1913&2114&7258
% & 389&626&749&1027&1539&1597
% & 47.0\%&36.0\%&38.8\%&29.0\%&41.0\%&23.9\%\\
 
CoCo'
& 790&1639&1875&1913&2114&7258
& 389&626&749&1027&1539&1597
& 51.5\%&41.7\%&73.2\%&40.6\%&43.7\%&40.2\%\\
 
 
\bvdesign-CoCo
& 758&1595&1919&1949&2145&7562
& 355&562&695&907&1320&1465
& 45.8\%&39.5\%&70.7\%&39.8\%&41.4\%&39.9\%\\
 
\csquared-CoCo
& 771&1637&1852&1900&2074&7237
& 358&560&613&898&1342&1448
& 45.8\%&35.7\%&38.5\%&28.9\%&40.7\%&23.9\%\\
 
\midrule
 
Marisa
& 234&470&676&712&521&1270
& 272&460&663&669&991&1098
& \cellcolor{blue!20}33.3\%&35.3\%&76.4\%&40.3\%&32.4\%&37.3\%\\
 
\bvdesign-Marisa
& 131&1046&1175&991&865&2031
& 178&\cellcolor{blue!20}310&443&\cellcolor{green!25}483&739&\cellcolor{blue!20}831
& 37.5\%&36.4\%&75.4\%&39.5\%&34.4\%&36.4\%\\
 
\csquared-Marisa
& 168&1064&1139&933&799&1613
& 182&\cellcolor{green!25}308&\cellcolor{blue!20}439&\cellcolor{blue!20}500&732&\cellcolor{green!25}829
& 37.5\%&30.9\%&41.0\%&27.6\%&32.9\%&20.7\%\\
 
Marisa-1
& 239&558&1002&945&599&1405
& 301&583&968&960&1244&1538
& \cellcolor{green!25}29.5\%&\cellcolor{green!25}22.0\%&34.3\%&25.2\%&\cellcolor{green!25}25.1\%&28.6\%\\
 
\bvdesign-Marisa-1
& 141&1229&2275&1476&1025&2209
& 178&425&686&746&967&1229
& 37.5\%&26.5\%&\cellcolor{blue!20}33.8\%&26.3\%&28.9\%&28.2\%\\
 
\csquared-Marisa-1
& 169&1294&2369&1473&997&1938
& 196&424&\cellcolor{green!25}436&764&979&1306
& 36.2\%&26.3\%&41.0\%&\cellcolor{blue!20}21.5\%&\cellcolor{blue!20}28.3\%&\cellcolor{blue!20}14.7\%\\
 
\midrule
 
PDT
& 624&1582&2947&4593&2040&3912
& 356&560&795&841&1022&1138
& 33.7\%&\cellcolor{blue!20}26.2\%&\cellcolor{green!25}28.4\%&\cellcolor{green!25}19.7\%&30.1\%&\cellcolor{green!25}13.0\%\\
 
ART
& \cellcolor{green!25}73&\cellcolor{green!25}83&\cellcolor{green!25}76&\cellcolor{green!25}115&\cellcolor{green!25}92&\cellcolor{green!25}289
& \cellcolor{green!25}132&347&533&728&\cellcolor{green!25}713&1611
& 346.8\%&156.2\%&115.0\%&59.1\%&165.7\%&27.6\%\\
 
C-ART
& 119&\cellcolor{blue!20}150&\cellcolor{blue!20}179&\cellcolor{blue!20}178&\cellcolor{blue!20}220&404
& \cellcolor{blue!20}148&401&674&750&810&1758
& 155.0\%&68.6\%&53.8\%&27.1\%&77.5\%&16.0\%\\
 
c-trie++
& 401&654&365&344&644&\cellcolor{blue!20}315
& 217&495&683&915&\cellcolor{blue!20}716&1507
& 1540.1\%&747.0\%&513.5\%&290.6\%&788.5\%&157.1\%\\
 
z-fast trie
& 447&779&989&1587&1168&2531
& 541&1158&1394&2016&1192&2838
& 1084.5\%&586.6\%&400.2\%&223.9\%&546.6\%&92.7\%\\
 
 
          \bottomrule
          \end{tabular}}
        }
    \end{center}
 
\label{tab:c2-full}
 
\end{table*}

\iffalse
\begin{table*}[t]

\caption{
%\todo{add more recursion levels for the other to get the pareto frontier?}
Build time,
query latency,
and space usage on the datasets in~\tabref{datasets}.
%Bolded numbers are the best in each column.
The prefixes
        \bvdesign/\csquared indicate that the
        \bvdesign/\bvdesign+\compressionscheme optimizations are applied.  We
        fix the tail container to be FSST with the \compressionscheme
        optimization, but we report results with the sorted tail container for
        the ablation study on the \bvdesign optimization. Marisa-1 means the
        Marisa trie with 1 recursion level.  For Marisa, we show all data points
        before \compressionscheme stops recursion.  For all other tries, the
        recursion level was set to 0 for a fair
        comparison.
}
    \begin{center}
        \newcolumntype{P}{>{\raggedleft\arraybackslash}p{.24in}}
        \resizebox{\textwidth}{!}{
\newrev{          \begin{tabular}{r P P P P P P P P P P P P P P P P P P}

            \toprule

             & \multicolumn{6}{c}{\textit{\buildtimeshort}   \textit{(ns/key)}}
             & \multicolumn{6}{c}{\textit{\queryshort} \textit{ (ns/key)}}
             & \multicolumn{6}{c}{\textit{\sizeshort} \textit{(\% of original size)}} \\

             \cmidrule(lr){2-7}
             \cmidrule(lr){8-13}
             \cmidrule(lr){14-19}

           \textit{Trie} &
           \textit{\wordsdata} & \textit{\urldata} & \textit{\dnadata} &
           \textit{\xmldata} & \textit{\wikidata} & \textit{\logdata} &
           \textit{\wordsdata} & \textit{\urldata} & \textit{\dnadata} &
           \textit{\xmldata} & \textit{\wikidata} & \textit{\logdata} &
           \textit{\wordsdata} & \textit{\urldata} & \textit{\dnadata} &
           \textit{\xmldata} & \textit{\wikidata} & \textit{\logdata} \\

            \midrule

FST
& 309&691&802&1157&716&2139
& 403&610&844&1541&1424&4003
& 45.5\%&35.5\%&74.3\%&39.2\%&37.9\%&39.5\%\\

\bvdesign-FST
& 90&1043&1148&897&762&2173
& 228&404&562&835&1242&2197
& 40.4\%&39.8\%&75.9\%&39.7\%&38.4\%&40.4\%\\

\csquared-FST
& 101&1068&1114&894&729&1821
& 227&404&561&842&1258&2154
& 40.4\%&36.0\%&43.8\%&28.8\%&37.7\%&24.4\%\\

\midrule

% CoCo'
% & 790&1639&1875&1913&2114&7258
% & 389&626&749&1027&1539&1597
% & 47.0\%&36.0\%&38.8\%&29.0\%&41.0\%&23.9\%\\

CoCo'
  & 790&1639&1875&1913&2114&7258
  & 389&626&749&1027&1539&1597
  & 51.5\%&41.7\%&73.2\%&40.6\%&43.7\%&40.2\%\\


\bvdesign-CoCo
& 758&1595&1919&1949&2145&7562
& 355&562&695&907&1320&1465
& 45.8\%&39.5\%&70.7\%&39.8\%&41.4\%&39.9\%\\

\csquared-CoCo
& 771&1637&1852&1900&2074&7237
& 358&560&613&898&1342&1448
& 45.8\%&35.7\%&38.5\%&28.9\%&40.7\%&23.9\%\\

\midrule

Marisa
& 234&470&676&712&521&1270
& 272&460&663&669&991&1098
& 33.3\%&35.3\%&76.4\%&40.3\%&32.4\%&37.3\%\\

\bvdesign-Marisa
& 131&1046&1175&991&865&2031
& 178&310&443&483&739&831
& 37.5\%&36.4\%&75.4\%&39.5\%&34.4\%&36.4\%\\

\csquared-Marisa
& 168&1064&1139&933&799&1613
& 182&308&439&500&732&829
& 37.5\%&30.9\%&41.0\%&27.6\%&32.9\%&20.7\%\\

Marisa-1
& 239&558&1002&945&599&1405
& 301&583&968&960&1244&1538
& 29.5\%&22.0\%&34.3\%&25.2\%&25.1\%&28.6\%\\

\bvdesign-Marisa-1
& 141&1229&2275&1476&1025&2209
& 178&425&686&746&967&1229
& 37.5\%&26.5\%&33.8\%&26.3\%&28.9\%&28.2\%\\

\csquared-Marisa-1
& 169&1294&2369&1473&997&1938
& 196&424&436&764&979&1306
& 36.2\%&26.3\%&41.0\%&21.5\%&28.3\%&14.7\%\\

\midrule

PDT
& 624&1582&2947&4593&2040&3912
& 356&560&795&841&1022&1138
& 33.7\%&26.2\%&28.4\%&19.7\%&30.1\%&13.0\%\\

ART
& 73&83&76&115&92&289
& 132&347&533&728&713&1611
& 346.8\%&156.2\%&115.0\%&59.1\%&165.7\%&27.6\%\\

C-ART
& 119&150&179&178&220&404
& 148&401&674&750&810&1758
& 155.0\%&68.6\%&53.8\%&27.1\%&77.5\%&16.0\%\\

c-trie++
  & 401&654&365&344&644&315
  & 217&495&683&915&716&1507
  & 1540.1\%&747.0\%&513.5\%&290.6\%&788.5\%&157.1\%\\

z-fast trie
  & 447&779&989&1587&1168&2531
  & 541&1158&1394&2016&1192&2838
  & 1084.5\%&586.6\%&400.2\%&223.9\%&546.6\%&92.7\%\\


          \bottomrule
          \end{tabular}}
        }
    \end{center}

\label{tab:c2-full}

\end{table*}
\fi



%%% Local Variables:
%%% mode: latex
%%% TeX-master: "../sample-sigconf"
%%% End:

\subsection{Ablation Study and Sensitivity Analysis}\label{sec:ablation-study}

\figref{intro-results-summary} shows ablation results of \bvdesign and \csquared
across configurations of FST, CoCo, and Marisa. \tabref{c2-full} contains the
full data for the different \csquared configurations discussed
in~\secreftwo{ablation-study}{performance-comparison}.

\newmymarginpar{Meta 5, R3 O4}\newnew{\para{Effect of hybrid
    \bv}~\tabref{fst-ablation-rel} evaluates the effect of switching to
  LOUDS-Sparse in FST~\cite{surf}. As mentioned in~\secref{preliminaries}, the
  original FST introduced a hybrid \bv with LOUDS-Sparse and LOUDS-Dense. We
  find that the hybrid \bv improves the performance of both build and query in
  the baseline FST (with no special tail container). However, the effect is much
  smaller on \csquared-FST with the FSST tail container, so we switch to the fully
  sparse version in the final \csquared configuration.}




%   \para{Effect of cache-optimized \bv design} First, we isolate the effect of
%   $C_1$, the cache-optimized \bv design. We perform the ablation study on the
%   datasets in~\tabref{datasets} (1) with the sorted tail container and no
%   recursion to avoid confounding factors from path compression, and (2) with all
%   branching labels stored out of place to suppress intra-node search effects.
% % We omit \texttt{wiki*} and \texttt{log*} from the table due to space
% % limitations and because the original CoCo-trie runs out of memory and cannot
% % build. However, as mentioned in~\secref{experimental-setup}, we introduce an
% % optimized build scheme CoCo' that addresses this issue. We observe similar
% % trends on these two larger datasets.

%   The ablation study in \tabref{c2-full} shows that the $C_1$ optimization
%   improves query performance by \bvfstspeedup, \bvcocospeedup, and \bvmarisaspeedup,
%   respectively, compared to the original FST, CoCo-trie, and Marisa.
%   %As shown
%   %in \tabref{cache-motivation}, \bvdesign reduces cache misses by an average of
%   %$1.51\times$ for \csquared-FST and $1.23\times$ for \csquared-Marisa.

%   The bitvector redesign has the least impact in the CoCo-trie because the query
%   time in CoCo is dominated by the cost to read the encodings. However, FST and
%   Marisa improve significantly because \bv cache misses dominate query time.

%   \begin{table}[t]
  \caption{Bitvector Operation Latency (ns) on XML Dataset (2.1M keys).
  \textit{Speedup} = Baseline / \bvdesign. Bold entries indicate where \bvdesign
  outperforms the baseline.}
  \centering
  \small
 \newrev{ \begin{tabular}{l l r r r}
  \toprule
  \textit{Trie} & \textit{Operation} & \textit{Baseline (ns)} & \textit{\bvdesign (ns)} & \textit{Speedup}
  \\
  \midrule
  \multirow{4}{*}{FST}
    & \texttt{get}        & 1.48  & 1.52           & $0.98\times$ \\
    & \texttt{leaf\_id}   & 10.32 & \textbf{1.45}  & $\mathbf{7.13\times}$ \\
    & \texttt{degree}     & 1.09  & 2.75           & $0.40\times$ \\
    & \texttt{child\_pos} & 48.62 & \textbf{17.09} & $\mathbf{2.84\times}$ \\
  \midrule
  \multirow{5}{*}{CoCo}
    & \texttt{get}           & 0.64  & 1.06           & $0.60\times$ \\
    & \texttt{leaf\_id}      & 2.10  & \textbf{1.72}  & $\mathbf{1.23\times}$ \\
    & \texttt{internal\_id}  & 4.18  & \textbf{2.82}  & $\mathbf{1.48\times}$ \\
    & \texttt{degree}        & 1.83  & 2.34           & $0.78\times$ \\
    & \texttt{child\_pos}    & 22.60 & \textbf{17.87} & $\mathbf{1.26\times}$ \\
  \midrule
  \multirow{5}{*}{Marisa}
    & \texttt{get}         & 0.98  & 1.48           & $0.66\times$ \\
    & \texttt{link\_id}    & 8.20  & \textbf{1.23}  & $\mathbf{6.66\times}$ \\
    & \texttt{leaf\_id}    & 9.00  & \textbf{1.38}  & $\mathbf{6.53\times}$ \\
    & \texttt{child\_pos}  & 30.31 & \textbf{16.62} & $\mathbf{1.82\times}$ \\
    & \texttt{parent\_pos} & 28.68 & \textbf{16.29} & $\mathbf{1.76\times}$ \\
  \bottomrule
  \end{tabular}}
  \label{tab:trie-op}
\end{table}

% \tabref{trie-op} reports the effect of \bvdesign on the fine-grained \bv
% operations that make up succinct-trie navigation.\mymarginpar{R1 D5} The results
% demonstrate that \bvdesign improves the performance of key \newrev{rank-based
%   operations such as \textsc{leaf\_id} and \textsc{internal\_id} by
%   $1.2-7.1\times$.} We observe similar trends for trie-specific operations, such
% as parent navigation in Marisa. \newrev{The only operators where the baseline is
%   faster are \textsc{get} and \textsc{degree}, which perform linear scans
%   through bits and do not use rank and select. However, neither of these lie on
%   the critical path of a full lookup or range query}.




\para{Effect of cache-optimized \bv design}
  First, we isolate the effect of \bvdesign, the cache-optimized \bv design.
  We perform the ablation study on the datasets in~\tabref{datasets} (1) with the
  sorted tail container and no recursion to avoid confounding factors from path
  compression, and (2) with all branching labels stored out of place to suppress
  intra-node search effects\newrev{, so that the remaining query-time difference
  is attributable to the \bv layout alone}.

  The ablation study in \tabref{c2-full} shows that the \bvdesign optimization
  improves query performance by \bvfstspeedup, \bvcocospeedup, and
  \bvmarisaspeedup, respectively, compared to the original FST, CoCo-trie, and
  Marisa. \newrev{These trie-level gains compound directly from the per-operator
  speedups in \tabref{trie-op}, since every full lookup issues
  $O(\mathit{depth})$ rank queries on the topology bitvector and a small
  constant number of select queries, each of which benefits from the co-located
  block metadata in our layout.}

  The bitvector redesign has the least impact in the CoCo-trie because the query
  time in CoCo is dominated by the cost to read the encodings, and each
  CoCo macro-node already amortises several traversal steps over one \bv
  access. However, FST and Marisa improve significantly because \bv cache
  misses dominate query time: each rank query in their baselines incurs
  two cache misses --- one for the block summary and one for the bit field ---
  which our co-located layout collapses to a single line.
  %eliminating the dominant component of the per-node \bv cost}.

  \begin{table}[t]
  \caption{Bitvector Operation Latency (ns) on XML Dataset (2.1M keys).
  \textit{Speedup} = Baseline / \bvdesign. Bold entries indicate where \bvdesign
  outperforms the baseline.}
  \centering
  \small
 \newrev{ \begin{tabular}{l l r r r}
  \toprule
  \textit{Trie} & \textit{Operation} & \textit{Baseline (ns)} & \textit{\bvdesign (ns)} & \textit{Speedup}
  \\
  \midrule
  \multirow{4}{*}{FST}
    & \texttt{get}        & 1.48  & 1.52           & $0.98\times$ \\
    & \texttt{leaf\_id}   & 10.32 & \textbf{1.45}  & $\mathbf{7.13\times}$ \\
    & \texttt{degree}     & 1.09  & 2.75           & $0.40\times$ \\
    & \texttt{child\_pos} & 48.62 & \textbf{17.09} & $\mathbf{2.84\times}$ \\
  \midrule
  \multirow{5}{*}{CoCo}
    & \texttt{get}           & 0.64  & 1.06           & $0.60\times$ \\
    & \texttt{leaf\_id}      & 2.10  & \textbf{1.72}  & $\mathbf{1.23\times}$ \\
    & \texttt{internal\_id}  & 4.18  & \textbf{2.82}  & $\mathbf{1.48\times}$ \\
    & \texttt{degree}        & 1.83  & 2.34           & $0.78\times$ \\
    & \texttt{child\_pos}    & 22.60 & \textbf{17.87} & $\mathbf{1.26\times}$ \\
  \midrule
  \multirow{5}{*}{Marisa}
    & \texttt{get}         & 0.98  & 1.48           & $0.66\times$ \\
    & \texttt{link\_id}    & 8.20  & \textbf{1.23}  & $\mathbf{6.66\times}$ \\
    & \texttt{leaf\_id}    & 9.00  & \textbf{1.38}  & $\mathbf{6.53\times}$ \\
    & \texttt{child\_pos}  & 30.31 & \textbf{16.62} & $\mathbf{1.82\times}$ \\
    & \texttt{parent\_pos} & 28.68 & \textbf{16.29} & $\mathbf{1.76\times}$ \\
  \bottomrule
  \end{tabular}}
  \label{tab:trie-op}
\end{table}

  \tabref{trie-op} reports the effect of \bvdesign on the fine-grained \bv
  operations that make up succinct-trie navigation. The results demonstrate that
  \bvdesign improves the performance of key \newrev{rank-based operations such
  as \texttt{leaf\_id} and \texttt{internal\_id} by $1.2-7.1\times$.}
  We observe similar trends for trie-specific operations, such as parent
  navigation in Marisa\newrev{ ($1.76\times$ on \texttt{parent\_pos}) and child
  navigation across all three tries ($1.3-2.8\times$ on \texttt{child\_pos})}.
  \newrev{The only operators where the baseline is faster are \texttt{get} and
  \texttt{degree}, which perform linear scans through bits and do not use rank
  and select\newrev{; interleaving block metadata with the bit field slightly
  reduces the effective bit density per cache line on such scans}. However,
  neither of these lie on the critical path of a full lookup or range
  query\newrev{, and both are invoked $O(1)$ times per query at sub-3-ns cost
  each, so the few nanoseconds lost are dwarfed by the tens of nanoseconds
  saved on every rank query}}.




\begin{figure*}[t]
    \centering
    \includegraphics[width=\textwidth]{figs/recursion_pareto_row.pdf}
    \caption{
    Space-latency Pareto frontier under recursion depths
    $\rho\in\{0,1,2\}$ for the three trie variants across all six datasets.}
    \label{fig:recursion-pareto}
\end{figure*}






\para{Effect of unary-path compression} \new{Next, we
  evaluate the effect of the \compressionscheme compression
  scheme. Specifically, for all tries, we measure the effect of replacing the
  sorted tail container with \fsst. Furthermore, we test different recursion
  levels to evaluate the space-time tradeoff exposed by more recursions.}
  %.~\figref{tuning-results}
 % visualizes the results of this sensitivity analysis on the \texttt{log}
 % dataset, and ~\tabref{c2-full} contains the full data.}

\newrev{Replacing the sorted tail container with \fsst improves space usage by
  \csqfstspacesaving, \csqcocospacesaving, and \csqmarisaspacesaving over
  \bvdesign-FST, \bvdesign-CoCo, and \bvdesign-Marisa, respectively. The effect
  on query performance is minimal (1-1.05$\times$ on average).
}






     % Fewer recursion levels expose Pareto-optimal space-time trade-offs.
     % For example, on \csquared-Marisa,
     %  \newmymarginpar{R3 D9} \newrev{on the \logdata dataset, no recursion takes 829 ns/query, while one level of recursion takes 1308 ns/query.
     %  On the other hand, the space usage improves from 20.7\% to 14.7\%. Similarly, on the \wikidata dataset, no recursion takes 732 ns/query compared to 979 ns/query with one level of recursion.
     %  For \wikidata, one level of recursion improves the space usage from 32.9\% to 28.3\%.}

     %     Beyond this point, further recursion's space savings do not justify the
     % performance penalty. For example, if we continue to recurse \csquared-Marisa-1
     % on \xmldata/\logdata, %\todo{update}
     % the relative space usages improve by about \newrev{3$\times$} (which
     % corresponds to a 15-20\% improvement relative to the original data size), but
     % the query latencies degrade by about \newrev{1.5$\times$}. We observe similar
     % results for the other tested tries and will include the relevant data in the
     % \newmymarginpar{R3 D7}\newrev{arXiv} version.\todo{put the actual data in?}


As shown in~\figref{recursion-pareto}, recursion levels 0 and 1 expose Pareto-optimal space-time trade-offs.
  For example, on \csquared-Marisa, on the \logdata dataset, no recursion takes
  829 ns/query, while one level of recursion takes 1308 ns/query. On the other
  hand, the space usage improves from 20.7\% to 14.7\%. Similarly, on the
  \wikidata dataset, no recursion takes 732 ns/query compared to 979 ns/query
  with one level of recursion. For \wikidata, one level of recursion improves
  the space usage from 32.9\% to 28.3\%.

  The same pattern holds for \csquared-FST and \csquared-CoCo. On the \logdata
  dataset, \csquared-FST without recursion takes 2187 ns/query at 24.4\% space,
  while one level of recursion takes 2502 ns/query at 20.1\% space; on the same
  dataset, \csquared-CoCo without recursion takes 1463 ns/query at 23.9\% space,
  while one level of recursion takes 1745 ns/query at 19.6\% space. The same
  trade also holds on \wikidata, \xmldata, \urldata, and \dnadata: across our
  six datasets, the absolute space savings from recursion levels 0 to 1 range
  from 0 percentage points (e.g., \wordsdata, whose avg.\ longest common prefix
  of 6 leaves no recursive sub-prefix structure to extract) to roughly 6
  percentage points (\logdata across all three variants), while the
  corresponding latency degradation stays below 1.5$\times$ in every cell.

  Beyond one level of recursion, further recursion's space savings rarely justify the
  performance penalty. For example, if we continue to recurse \csquared-Marisa-1
  on \logdata, the relative space usage improves by 1.6$\times$ (which
  corresponds to a 5.3 percentage point improvement relative to the original
  data size, from 14.7\% to 9.4\%), but the query latency degrades by 1.25$\times$
  (from 1308 to 1593 ns/query); on \xmldata, the same step yields no measurable
  further space improvement. In fact, \logdata is the only dataset on which
  $\rho{=}2$ strictly Pareto-dominates $\rho{=}1$ across all three tries, owing
  to its 113-character average key length harbouring nested common prefixes
  that a second recursion round can still extract. We observe one further
  anomaly: on \dnadata, \csquared-Marisa is fully insensitive to recursion ---
  its size remains at 41.1\% of the original data and its latency varies by less
  than 1\% across $\rho \in \{0,1,2\}$ --- because its unary-path encoder
  already absorbs DNA's unique-suffix structure at $\rho{=}0$, leaving no
  inter-tail prefix redundancy for the adaptive cost-based recursive
  string-pool builder to exploit.

  Going forward, we limit the recursion level of \csquared-FST and \csquared-CoCo
  to 0 to match the original designs. For \csquared-Marisa, we use the recursion
  level prescribed by the adaptive compression scheme. We expose recursion as
  an option for space-constrained settings: dialling \csquared-Marisa up to
  $\rho{=}2$ on long-key datasets like \logdata trades a further 5 percentage
  points of space for a 1.25$\times$ query latency penalty, while on every
  other dataset the same dial yields at most one percentage point of additional
  space and is therefore Pareto-dominated by $\rho{=}1$.












\begin{figure*}[t]
    \centering

    \begin{subfigure}[t]{0.32\textwidth}
        \centering
        \includegraphics[width=\linewidth]{figs/range_latency_fst_vs_c2fst_k1.pdf}
        % \caption{$w=10$}
        \label{fig:range-fst-w10}
    \end{subfigure}
    \hfill
    \begin{subfigure}[t]{0.32\textwidth}
        \centering
        \includegraphics[width=\linewidth]{figs/range_latency_fst_vs_c2fst_k100.pdf}
        % {figs/range_latency_fst_vs_c2fst_w100.pdf}
        % \caption{$w=100$}
        \label{fig:range-fst-w100}
    \end{subfigure}
    \hfill
    \begin{subfigure}[t]{0.32\textwidth}
      \centering
      \includegraphics[width=\linewidth]{figs/range_latency_fst_vs_c2fst_k10000.pdf}
        % {figs/range_latency_fst_vs_c2fst_w1000.pdf}
        % \caption{$w=1000$}
        \label{fig:range-fst-w1000}
    \end{subfigure}
    \caption{ \newnew{ Range-query latency comparison between FST-baseline and
        \csquared-FST across six datasets under different range widths $k$.  Each
        subplot corresponds to a fixed range width.  Bars report absolute
        latency (left axis), while the red line indicates the speedup of
        \csquared-FST over the baseline (right axis).} }

    \label{fig:range-fst-vs-c2fst}
\end{figure*}





\subsection{Comparison to State-of-the-art
  Tries}\label{sec:performance-comparison}\mymarginpar{R2 D10}

We compare \csquared tries with all baselines.\rev{~\tabref{c2-full} reports the
  full data and the configurations prescribed by the \compressionscheme, the
  adaptive compression algorithm.}

\para{Existence queries} \rev{Among all tries, ART always has the best query
  performance at the cost of high space usage. Child navigation in ART simply
  requires following a pointer, while succinct tries must perform complex \bv
  operations. \mymarginpar{R1 D5, R3 D2}\new{Furthermore,
    C-ART~\cite{hybrid_index} uses an internal-node organization that trades
    search performance for space savings over ART.}} \newmymarginpar{Meta 4, R2
  O2}\newrev{Both c-trie++ and z-fast trie are slower than ART for queries while
  consuming significantly more space.}


\new{\tabref{c2-full} demonstrates that \compressionscheme does not affect the
  query-time gains from \bvdesign across tries and datasets.}  \rev{Among the
  succinct tries}, \csquared-Marisa achieves the lowest query latency with a
space usage within $1.2\times$ of the best. On average, \csquared-Marisa
improves the query latency of the original Marisa by \csqmarisaspeedup.

On \texttt{words} and \texttt{url}, \csquared-Marisa-1's space consumption is
relatively high compared to \rev{Marisa}-1 \mymarginpar{R3 O4}\rev{because in
  this case, the \compressionscheme scheme stores short unary paths in place
  rather than in the tail container (see~\secref{unary-path}). For example, on
  the \texttt{words} dataset, \csquared-Marisa contains only about .3$\times$ as many
  links as Marisa.}

% \todo{make sure these hold with the changes in fixing recursion
  % level}
\csquared-FST improves query latency/space usage by
\csqfstspeedup/\csqfstspacesaving on average compared with FST.
% while using similar or less memory (between $.73\times-1.12\times$).
\csquared-CoCo does not benefit as \rev{much} from \bvdesign as it incurs fewer
\bv operations, but still improves both space usage and query latency compared
to \rev{CoCo'}.

\para{Time-space tradeoff} \new{We next examine query time and space
  usage. \compressionscheme improves space usage by $1.34\times$ \rev{on
    average} compared to their original versions.}

\new{On average, with no recursion, switching from the
  sorted tail container to FSST reduces the space usage by \toplevelspacesavings at the
  cost of a $1.05\times$ query performance penalty and $1.1\times$ higher build
  time across all tested succinct tries.

  Additionally, for Marisa, on the datasets where \compressionscheme prescribes
  one level of recursion, the recursion improves the space efficiency by
  % \todo{update numbers to reflect marisa only}
  $1.29\times$ on average, but also incurs a query performance penalty of
  $1.43\times$ compared to no recursions.}  \rev{Specifically, recursion turns
  out to be extremely effective on the highly redundant \logdata dataset. For
  \csquared-Marisa, one level of recursion reduces the space cost by
  $1.74\times$, but increases the query time by $1.7\times$ compared to no
  recursions due to the presence of many recursive paths.}

\new{On \wordsdata, \urldata and \wikidata, \csquared-Marisa-1 has higher space
  usage (by at most $1.19\times$) than Marisa-1 because of the space overhead of
  storing the branching labels in place
  (see~\secref{unary-path}). %outweighs the savings of \fsst. %the FSST tail container in these cases.
}

On the \rev{larger datasets (e.g., \logdata)}, \csquared-FST and \csquared-CoCo
are generally less competitive than \csquared-Marisa in both query performance
and compression ratio, due to their handling of internal unary
paths. \csquared-FST does not handle such paths at all, \rev{and} \csquared-CoCo
does not address the data redundancy in such paths. Furthermore, unary paths
that do not fit the machine word size will be fragmented in \csquared-CoCo,
impairing data locality.

\newrev{\para{Range queries} \newmymarginpar{Meta 3, R1 O3}
  ~\figref{range-fst-vs-c2fst} shows that the \csquared optimizations speed up
  range queries by $1.2-1.5\times$ in FST, which was the only tested trie that
  natively supports successor queries. The locality improvements in
  \csquared-FST carry over into range queries, where each \texttt{next()} step
  can be resolved from data in the packed 256-bit blocks, which are already in
  cache. In contrast, the original FST fetches the equivalent information from
  three separately heap-allocated arrays, incurring up to three cache-line
  misses per step.

  The relative advantage of \csquared-FST increases monotonically with range
  width, with speedups on all but one configuration (\wikidata with $k =
  1$). The \wikidata contains many diverse suffixes, so the FSST tail container
  in \csquared-FST finds almost nothing to compress, as shown
  in~\tabref{c2-full}. However, \csquared-FST must incur additional indirections
  during traversal to check the \texttt{is\_link} array and fetch the compressed
  suffix, while the original FST can do a direct comparison.  }

\para{Build time}\label{sec:build-performance} While modern static succinct
tries excel at space efficiency and query latency, their build \rev{times are}
significantly worse (\rev{between 3.2-24.7$\times$}) than \rev{the build times}
of pointer-based tries such as ART and C-ART, especially \rev{on large
  datasets}.\newmymarginpar{R3 D8} \newrev{Depending on the dataset, the
  \csquared optimizations can increase the build time by about $2\times$ due to
  the time it takes to compress the tail container. We chose FSST as the tail
  container to try to minimize build times compared to approximate \repair,
  which can take another $2\times$ longer, as shown in PDT's build times.}
\rev{Furthermore, building is a one-time cost for succinct tries and can be
  amortized over many queries.}
  %Additionally, we restrict \compressionscheme
  %to the \fsst tail container for faster build times with minimal space cost
  %compared to approximate \repair.}
  Future work will parallelize the build phase of succinct tries.



%%% Local Variables:
%%% mode: latex
%%% TeX-master: "main"
%%% End:
